\section{Is the gap overfitting, or is it drift?}
\label{sec:split}

Section~\ref{sec:results} reports that both families are markedly worse on
recordings they never saw, and that the two differ in how much. That result is
open to an objection which this section closes. The objection is worth stating
in full, because it applies to any result obtained the way these were.

The trainer holds back the \emph{last} recordings of a session, not a random
selection. A single session drifts: the lighting changes as the afternoon goes
on, the garment is replaced slightly differently between takes, the tactile gels
pick up residue and wear. Every one of those changes accumulates towards the end
of the session, which is precisely where the held-back recordings are. A policy
could therefore be perfectly good at generalising and still score badly on them,
simply because they differ from everything it was shown. Read that way, the gap
is not a statement about the policy at all.

The control is a second retraining of each family on a \emph{random} seven held
back instead of the last seven, with the same architecture, the same optimisation
budget, the same seed and the same scoring. Table~\ref{tab:split} gives the
result and Figure~\ref{fig:split} draws it.

\begin{table}[h]
  \centering
  \caption{The generalisation gap under each hold-out. The gap is held-out error
  divided by trained-on error, so it is a ratio within one policy and does not
  depend on the differing planning horizons.}
  \label{tab:split}
  \small
  \begin{tabular}{lrrr}
    \toprule
    & \multicolumn{2}{c}{gap} & \\
    \cmidrule(lr){2-3}
    family & last seven & random seven & change \\
    \midrule
    action-chunking & 5.3 & 8.9 & $+68\,\%$ \\
    diffusion & 1.9 & 2.6 & $+37\,\%$ \\
    \bottomrule
  \end{tabular}
\end{table}

\begin{figure}[t]
  \centering
  \includegraphics[width=0.8\textwidth]{split_comparison.png}
  \caption{The generalisation gap under a temporal and a random hold-out. The
  dotted line marks a policy that does equally well on both halves. Colour
  follows the family, so a reader compares the pair within a colour rather than
  the heights across colours.}
  \label{fig:split}
\end{figure}

The direction settles the question. A random hold-out asks a policy to
interpolate: the recordings it must handle sit between recordings it was trained
on, so whatever drifted across the session drifted past them too and it has seen
both sides. The temporal hold-out asks it to extrapolate past the end of
everything it has seen. If drift were what produced the gap, the interpolating
case would be the easier of the two. It is the harder one, for both families and
by a wide margin. The degradation is therefore a property of recordings the
policy has not seen, and not of when in the session they were recorded.

One caveat survives, and it bounds the magnitudes rather than the direction. The
two hold-outs are different recordings, so the two ratios in
Table~\ref{tab:split} are not two measurements of a single quantity: some part
of the distance between them is that one set of seven recordings may simply be
harder than the other. Figure~\ref{fig:contact} shows at least one visibly
atypical recording among the random seven. The direction does not depend on
this, because it rests on a comparison the confound cannot reverse --- no
difference in difficulty between the two sets makes interpolation harder than
extrapolation as a matter of principle.

What the section does not establish is why the gap is as large as it is. Both
families were trained on fifty-eight recordings of a single task in a single
session, and the result says that is not enough for either of them to generalise
to a fifty-ninth. It does not say how many would be, nor whether variety across
sessions would help more than volume within one.
